\documentclass[11pt,a4paper]{article}

\usepackage[margin=1in]{geometry}
\usepackage{amsmath,amssymb}
\usepackage{physics}
\usepackage{booktabs}
\usepackage{hyperref}
\usepackage{enumitem}

\newenvironment{resultbox}[1][Result]{%
  \begin{quote}\noindent\textbf{#1.}\ }{\end{quote}}
\newenvironment{gapbox}[1][Open problem]{%
  \begin{quote}\noindent\textbf{#1.}\ }{\end{quote}}

\newcommand{\half}{\tfrac{1}{2}}
\newcommand{\mzero}{m_0}
\newcommand{\seriesname}{Saturated Superfluid Vacuum}
\newcommand{\papernumber}{Alpha}

\title{\textbf{\seriesname~\papernumber}\\[0.4em]
The Fine-Structure Constant as a Vortex Aspect Ratio}
\author{Stig Norland \\ Independent Researcher \\ Bergen, Norway}
\date{\today}

\begin{document}
\maketitle

\begin{abstract}
The fine-structure constant $\alpha \approx 1/137$ enters ordinary
quantum electrodynamics as the dimensionless strength of electromagnetic
coupling.  In the Saturated Superfluid Vacuum (SSV) programme it appears
more widely: as the strength of the chiral-shear sector, the ratio between
the electron vortex-ring radius and the healing length, the scale
$\mu_0=m_e/\alpha$ of the mass ladder, the small parameter controlling
finite-ring corrections, and the geometric lever in weak-sector
reconnection estimates.  This paper isolates the common content of these
appearances.

We do not claim here to derive the numerical value of $\alpha$.  Instead
we show that, within SSV, the problem of deriving $\alpha$ is equivalent to
a sharply defined defect-geometry problem: compute the equilibrium aspect
ratio $R^*/\xi$ of the stable toroidal electron vortex from the full
LogSE plus chiral-shear functional, without inserting $\alpha$ as an
empirical input, and show that
\[
        \frac{R^*}{\xi} = \alpha^{-1} \approx 137.036 .
\]
This reframes the fine-structure constant as a dimensionless eigenvalue of
the vacuum defect spectrum rather than as a fundamental coupling constant
postulated independently of the medium.

The paper also clarifies what this proposal is not.  Following the
propagation-speed audit in Paper~II, $\alpha$ should not be interpreted as
the ratio of the observed photon speed to the medium sound speed.  The
observed photon is instead expected to belong to the Goldstone/Bogoliubov
channel propagating at $c$, while $\alpha$ belongs to the static
chiral-shear sector that fixes Coulomb strength and defect geometry.  The
central open calculation is therefore not a wave-speed calculation but a
three-dimensional toroidal-vortex minimisation.
\end{abstract}

\newpage
\tableofcontents
\newpage

% ---------------------------------------------------------------------
\section{Introduction}
\label{sec:introduction}
% ---------------------------------------------------------------------

The fine-structure constant is one of the sharpest dimensionless numbers
in physics:
\begin{equation}
    \alpha = \frac{e^2}{4\pi\epsilon_0\hbar c}
    \approx 7.29735256\times 10^{-3},
    \qquad
    \alpha^{-1}\approx 137.036 .
    \label{eq:alpha_qed}
\end{equation}
In the Standard Model it is a measured low-energy coupling, whose running
with scale is computed once the empirical value is fixed.  Its numerical
value is not derived from the structure of the theory.

The SSV programme changes the location of the question.  In Paper~I
\cite{SSV-I}, particles were identified with topological defects of a
saturated logarithmic superfluid: leptons as toroidal vortex rings and
baryons as breather modes stabilised by knotted vortex skeletons.  In
Paper~II \cite{SSV-II}, the force sector was developed and the earlier
wave-speed interpretation of $\alpha$ was refined: the observed photon
should not be the chiral-shear mode propagating at $c_\perp$, but rather a
Goldstone/Bogoliubov-channel excitation propagating at $c$.  The
chiral-shear sector remains physically real, but its role is the static
near-field Coulomb interaction and the geometry of charged defects.

This paper extracts the resulting implication:
\begin{quote}
\emph{In SSV, deriving the fine-structure constant is equivalent to
deriving the aspect ratio of the stable toroidal electron vortex.}
\end{quote}

The key target is
\begin{equation}
    R^* = \frac{\xi}{\alpha},
    \qquad
    \frac{R^*}{\xi} = \alpha^{-1}.
    \label{eq:target_ratio}
\end{equation}
Here $\xi$ is the electron healing length, the minor/core scale of the
vortex ring, and $R^*$ is the equilibrium major radius of the electron
torus.  Paper~I used $\alpha$ as an empirical input to obtain
$R^*/\xi\approx137$.  The present paper reverses the logic: the first
principles calculation that remains to be done is the minimisation that
outputs this number.

\begin{resultbox}[Claim of this paper]
The SSV framework does not yet derive $\alpha$.  It reduces the problem to
a concrete eigenvalue calculation: solve the full 3D LogSE plus
chiral-shear toroidal-vortex variational problem and compute the
dimensionless equilibrium aspect ratio $R^*/\xi$.  If the result is
$137.036\ldots$, the fine-structure constant is geometrically derived as
$\alpha=\xi/R^*$.
\end{resultbox}

% ---------------------------------------------------------------------
\section{Inputs from the SSV Series}
\label{sec:inputs}
% ---------------------------------------------------------------------

Only three pieces of the previous papers are needed here.

\paragraph{The saturated medium.}
The vacuum is described by a macroscopic complex order parameter
$\Psi(\mathbf{x},t)$.  Its Madelung form gives a density and phase,
\[
    \Psi = \sqrt{\rho}\,e^{i\theta},
    \qquad
    \mathbf{v} = \frac{\hbar}{\mzero}\nabla\theta ,
\]
with quantised circulation
\begin{equation}
    \oint \mathbf{v}\cdot d\mathbf{l}
    = n\kappa_0,
    \qquad
    \kappa_0 = \frac{h}{\mzero}.
    \label{eq:circulation}
\end{equation}
The logarithmic equation of state fixes the sound speed at saturation and
the healing length
\begin{equation}
    c_s(\rho_0)=c,
    \qquad
    \xi = \frac{\hbar}{\mzero c}.
    \label{eq:xi}
\end{equation}

\paragraph{The charged toroidal defect.}
The electron is identified with a stable toroidal vortex ring.  Its
minor/core radius is set by $\xi$; its major radius $R$ is a variational
degree of freedom.  Charge is the chirality of the toroidal winding: the
two signs of electric charge correspond to the two orientations of the
vortex-ring phase winding.

\paragraph{The chiral-shear sector.}
Paper~I introduced a chiral-shear contribution to the functional in order
to make transverse, handed flow gradients energetic.  Paper~II clarified
that the empirical $\alpha$ in this sector should be read as a
dimensionless chiral-shear coupling, not as the speed of the observed
photon.  The Coulomb force is then recovered as a Bernoulli pressure of the
chiral-shear near field,
\begin{equation}
    F_C(r) = \frac{\alpha\hbar c}{r^2},
    \label{eq:coulomb_ssv}
\end{equation}
after matching the elementary charge through
$e^2/(4\pi\epsilon_0)=\alpha\hbar c$.

These inputs deliberately leave the numerical value of $\alpha$ open.
That is the point of the present paper.

% ---------------------------------------------------------------------
\section{Where \texorpdfstring{$\alpha$}{alpha} Appears}
\label{sec:appearances}
% ---------------------------------------------------------------------

The same dimensionless number appears in several places in the SSV series.
The table below collects the appearances that matter for the derivation
problem.

\begin{center}
\renewcommand{\arraystretch}{1.35}
\begin{tabular}{lll}
\toprule
\textbf{Sector} & \textbf{Appearance of $\alpha$} & \textbf{SSV reading} \\
\midrule
Electromagnetic near field
  & $F_C=\alpha\hbar c/r^2$
  & chiral-shear coupling strength \\
Electron geometry
  & $R^*/\xi=\alpha^{-1}$
  & toroidal aspect ratio \\
Mass ladder
  & $\mu_0=m_e/\alpha$
  & flux-tube energy scale \\
Finite ring corrections
  & $(\xi/R^*)^2=\alpha^2$
  & small expansion parameter \\
Weak reconnection estimates
  & $E_{\rm recon}\sim m_e/\alpha^2$
  & ring-scale cavity energy \\
Gravity consistency
  & $G\propto \alpha^2/(N_p^2m_e^2)$
  & mass-ladder bookkeeping \\
\bottomrule
\end{tabular}
\end{center}

This pattern is suggestive but not itself a derivation.  The danger is to
mistake recurrence for explanation.  A number that appears everywhere in a
framework may be fundamental, or it may be the shadow of one deeper
geometric quantity.  The SSV proposal is the latter: all appearances of
$\alpha$ trace back to the relative scale of the electron torus and its
core.

\begin{resultbox}[Reduction]
If the full toroidal-vortex calculation yields $R^*/\xi=137.036\ldots$,
then the Coulomb coupling, the mass-ladder scale, and the small expansion
parameter all inherit the same number from one defect-geometry eigenvalue.
\end{resultbox}

% ---------------------------------------------------------------------
\section{What \texorpdfstring{$\alpha$}{alpha} Is Not}
\label{sec:not}
% ---------------------------------------------------------------------

The early SSV drafts treated $\alpha$ as a wave-speed ratio in the form
$c_\perp/c$ or $c_\perp^2/c^2$, depending on convention.  This was a useful
scaffold because it placed $\alpha$ in the medium rather than in a
separate gauge sector.  But the propagation-speed audit in Paper~II shows
that this cannot be the final physical interpretation of the observed
photon.  Light in vacuum propagates at $c$, and so does gravitational
radiation to the accuracy constrained by GW170817.  The channel that
supports physical radiation is therefore the Goldstone/Bogoliubov channel,
not the slower chiral-shear mode.

This distinction is not a retreat from the SSV interpretation of
$\alpha$; it sharpens it.  The role of $\alpha$ is not to slow photons.  It
is to set the strength and geometry of the chiral-shear sector:
\begin{equation}
    \alpha
    \quad\hbox{belongs to coupling and defect geometry, not to photon speed.}
    \label{eq:not_speed}
\end{equation}

The time-dilation check in Paper~II reinforces this split.  In a
density-depressed region near a mass, the longitudinal Bogoliubov speed
scales as
\begin{equation}
    \frac{c_s(\rho)}{c_s(\rho_0)}
    = \sqrt{\frac{\rho}{\rho_0}}
    \approx 1+\frac{\Phi}{c^2},
    \label{eq:time_dilation_alpha_paper}
\end{equation}
which has the correct weak-field sign for gravitational time dilation.  A
naive local-density reading of the chiral-shear speed has the opposite
sign.  The same audit therefore points to a clean separation:
\begin{align}
    c &: \hbox{Goldstone/Bogoliubov propagation, light, gravity, clocks},\\
    \alpha &: \hbox{chiral-shear coupling, Coulomb near field, defect aspect ratio}.
\end{align}

% ---------------------------------------------------------------------
\section{The Geometric Target}
\label{sec:target}
% ---------------------------------------------------------------------

The electron ring is characterised by a major radius $R$ and a minor/core
radius $a$.  The core minimisation is local and fixes
\begin{equation}
    a^* = \xi .
    \label{eq:core_target}
\end{equation}
The nontrivial calculation is the major-radius minimisation,
\begin{equation}
    R^* = \operatorname*{arg\,min}_{R}
    E_{\rm torus}[R,a^*;\,\mathcal{L}_{\rm LogSE}+\mathcal{L}_{\perp}] .
    \label{eq:R_min}
\end{equation}
The desired output is the dimensionless number
\begin{equation}
    r^* \equiv \frac{R^*}{\xi}.
    \label{eq:rstar}
\end{equation}
In Paper~I, the empirical input $\alpha$ was used to write
\[
    r^* = \alpha^{-1}.
\]
The alpha derivation problem is the inverse question:
\begin{equation}
    r^*_{\rm computed}
    \stackrel{?}{=}
    137.036\ldots .
    \label{eq:alpha_test}
\end{equation}

\begin{gapbox}[Open calculation]
The alpha problem is not solved by restating $R^*=\xi/\alpha$.  It is solved
only if the functional minimisation computes $R^*/\xi$ without inserting
$\alpha$ anywhere as an empirical coupling or target radius.
\end{gapbox}

% ---------------------------------------------------------------------
\section{Variational Form of the Problem}
\label{sec:variational}
% ---------------------------------------------------------------------

The most compact way to state the calculation is as a constrained
minimisation over toroidal vortex configurations:
\begin{equation}
    E[\Psi]
    =
    \int d^3x\,
    \left[
        \frac{\hbar^2}{2\mzero}|\nabla\Psi|^2
        + V_{\rm LogSE}(|\Psi|^2)
        + \mathcal{E}_{\perp}(\Psi,\nabla\Psi)
    \right],
    \label{eq:functional}
\end{equation}
subject to the topological constraints
\begin{equation}
    \oint_{\rm poloidal}\nabla\theta\cdot d\mathbf{l}=2\pi n,
    \qquad
    \oint_{\rm toroidal}\nabla\theta\cdot d\mathbf{l}=2\pi m,
    \label{eq:constraints}
\end{equation}
with the electron corresponding to the minimal charged toroidal sector.
The chiral-shear term $\mathcal{E}_{\perp}$ must be written in a form that
does not already contain the observed $\alpha$ as a numerical input.  Its
dimensionless coefficient should either be fixed by saturation/topology or
absorbed into the same eigenvalue problem.

In the thin-ring language, the energy can be schematically decomposed as
\begin{equation}
    E_{\rm torus}(R,a)
    =
    E_{\rm core}(a)
    + E_{\rm flow}(R,a)
    + E_{\rm shear}(R,a)
    + E_{\rm curvature}(R,a)
    + \cdots ,
    \label{eq:energy_decomp}
\end{equation}
where
\begin{itemize}[leftmargin=2em]
    \item $E_{\rm core}$ fixes the minor radius $a^*=\xi$;
    \item $E_{\rm flow}$ is the circulation energy of the vortex ring;
    \item $E_{\rm shear}$ is the chiral contribution that resists or favours expansion;
    \item $E_{\rm curvature}$ contains Kelvin, bending, and finite-core corrections.
\end{itemize}
The output of the calculation is not a mass in isolation, but the
dimensionless stationary point
\begin{equation}
    \frac{\partial E_{\rm torus}}{\partial R}=0,
    \qquad
    \frac{\partial^2 E_{\rm torus}}{\partial R^2}>0,
    \qquad
    r^*=\frac{R^*}{\xi}.
    \label{eq:stationary}
\end{equation}

% ---------------------------------------------------------------------
\section{Avoiding the Circular Derivation}
\label{sec:circularity}
% ---------------------------------------------------------------------

The most important discipline in an alpha paper is to avoid proving
$\alpha$ by hiding it in the assumptions.  The following routes are
therefore disallowed as derivations:

\begin{enumerate}[label=(\roman*)]
    \item inserting $\lambda=\alpha^2\rho_0/\mzero$ and then obtaining
          $R^*/\xi=\alpha^{-1}$;
    \item imposing $e^2/(4\pi\epsilon_0)=\alpha\hbar c$ before solving the
          toroidal defect;
    \item choosing a boundary condition equivalent to $R^*/\xi=137$;
    \item fitting a numerical coefficient in the shear term to recover the
          observed value.
\end{enumerate}

A valid derivation must instead have the following form:
\begin{enumerate}[label=(D\arabic*)]
    \item Write the LogSE plus chiral-shear functional with all
          dimensionless coefficients fixed by symmetry, saturation, or
          topology.
    \item Solve the minimal charged toroidal defect in three dimensions.
    \item Extract the equilibrium ratio $R^*/\xi$.
    \item Only then identify $\alpha=\xi/R^*$ and compare with CODATA.
\end{enumerate}

\begin{resultbox}[Success criterion]
The alpha derivation succeeds if a calculation that contains no empirical
$\alpha$ produces
\[
    R^*/\xi = 137.036\ldots
\]
with a controlled error estimate.  It fails if the result is not close, or
if the value can be recovered only by inserting an equivalent dimensionless
constant by hand.
\end{resultbox}

% ---------------------------------------------------------------------
\section{Numerical Programme}
\label{sec:numerics}
% ---------------------------------------------------------------------

The practical route is numerical before it is analytic.  The cleanest
calculation is a 3D constrained minimisation in the toroidal topological
sector:
\begin{enumerate}[label=(N\arabic*)]
    \item Choose a cylindrical or toroidal computational grid with adaptive
          refinement near the core.
    \item Initialise $\Psi$ with a minimal charged toroidal vortex of major
          radius $R$ and core radius near $\xi$.
    \item Evolve by imaginary time or gradient flow under the full
          LogSE plus chiral-shear functional.
    \item Enforce the topological winding constraints so the defect cannot
          unwind.
    \item Measure the relaxed $R^*/\xi$ and its dependence on grid size,
          box size, and finite-core treatment.
    \item Repeat across allowed chiral-shear functional variants to test
          whether $137$ is robust or model-dependent.
\end{enumerate}

The calculation has a clear falsification structure.  A robust equilibrium
near $137$ would be strong evidence that $\alpha$ is indeed a vortex aspect
ratio.  A robust equilibrium at some unrelated number would falsify the
specific SSV alpha proposal, even if other parts of the framework survived.

% ---------------------------------------------------------------------
\section{Consequences of a Successful Derivation}
\label{sec:consequences}
% ---------------------------------------------------------------------

If $R^*/\xi=\alpha^{-1}$ is derived, several SSV claims become less
independent than they currently appear.

\paragraph{Electromagnetic coupling.}
The Coulomb strength would be inherited from the geometry of the charged
defect:
\[
    e^2/(4\pi\epsilon_0)=\hbar c\,\xi/R^*.
\]
The electromagnetic coupling would no longer be a separately inserted
number.

\paragraph{The mass ladder.}
The scale $\mu_0=m_e/\alpha$ would become
\[
    \mu_0 = m_e\,\frac{R^*}{\xi},
\]
so the hadronic ladder would inherit its spacing from the electron-ring
aspect ratio.

\paragraph{Weak-sector estimates.}
Quantities scaling as $m_e/\alpha^2$ would become ring-geometry energies
scaling as
\[
    m_e\left(\frac{R^*}{\xi}\right)^2 .
\]
The weak sector would still require the reconnection-cap calculation, but
its small parameter would be geometrically anchored.

\paragraph{Perturbative corrections.}
Finite-ring corrections use
\[
    \varepsilon = (\xi/R^*)^2 = \alpha^2 .
\]
If $\alpha$ is geometric, the perturbation expansion is an expansion in
the inverse aspect ratio of the electron torus.

% ---------------------------------------------------------------------
\section{Discussion}
\label{sec:discussion}
% ---------------------------------------------------------------------

This paper is intentionally narrower than the force-sector unification of
Paper~II.  Its purpose is to identify one load-bearing calculation.  The
fine-structure constant is not explained by showing that it appears in many
places.  It is explained only if all of those appearances trace back to a
single dimensionless eigenvalue of the vacuum defect spectrum.

The connectedness of the SSV series is therefore not a weakness here.  It
is the reason the alpha problem is worth isolating.  Coulomb strength,
electron geometry, mass-ladder spacing, weak reconnection scales, and
finite-ring corrections all point to the same ratio.  If the ratio can be
computed, the framework gains a central derivation.  If it cannot, the
failure is equally informative.

There is also an important conceptual payoff.  The propagation-speed audit
separates two roles that were previously entangled:
\[
    \hbox{propagation speed} \neq \hbox{coupling strength}.
\]
The observed photon belongs to the propagation channel at $c$.  The
fine-structure constant belongs to the defect and near-field channel.  This
makes the alpha problem more geometric, not less.

% ---------------------------------------------------------------------
\section{Conclusion}
\label{sec:conclusion}
% ---------------------------------------------------------------------

The SSV framework does not yet derive the fine-structure constant.  It
does, however, sharpen the question.  The fine-structure constant is not
treated as an unexplained electromagnetic input, nor as a photon-speed
ratio.  It is the dimensionless chiral-shear/topological scale that fixes
the electron torus relative to its healing length.

The derivation target is therefore:
\[
    \boxed{\quad R^*/\xi = \alpha^{-1} \approx 137.036 \quad}
\]
from the full 3D LogSE plus chiral-shear toroidal-vortex minimisation.
That calculation is small enough to be stated precisely and large enough to
matter.  It is the point at which the SSV programme either turns the
fine-structure constant into geometry, or learns exactly where the proposed
geometry fails.

\begin{thebibliography}{99}

\bibitem{SSV-I}
S.~Norland,
``Saturated Superfluid Vacuum I: Particle Defects and the
$\alpha$-Harmonic Mass Ladder,''
manuscript in preparation.

\bibitem{SSV-II}
S.~Norland,
``Saturated Superfluid Vacuum II: The Force Sector,''
manuscript in preparation.

\bibitem{Volovik2003}
G.~E.~Volovik,
\textit{The Universe in a Helium Droplet},
Oxford University Press, Oxford (2003).

\bibitem{Zloshchastiev2020}
K.~G.~Zloshchastiev,
``Superfluid vacuum theory and deformed dispersion relations,''
\textit{Int.~J.~Mod.~Phys.~A} \textbf{35}, 2040032 (2020);
arXiv:2011.11897.

\end{thebibliography}

\end{document}
